% Laget automatisk av elevark/lag_elevark.py fra sim/elektrostatisk-pingpong.html. Endringer her blir overskrevet.
\documentclass[11pt]{article}
\usepackage{elevark}
\begin{document}
\elevarkhode{Fysikk 2 · Elektriske felt og krefter}{Elektrostatisk ping-pong}{https://fysikklab.no/sim/elektrostatisk-pingpong.html}{Simuleringen}

En lett kule med ledende overflate henger i en tynn tråd mellom to plater med høy spenning. Når kula berører den positive plata, blir den positiv og skyves over til den negative. Der blir den negativ og skyves tilbake. Kula frakter ladning fram og tilbake, og det blir en liten strøm.

\section*{Oppgaver}

Start med $U = 10\,\mathrm{kV}$, $d = 8{,}0\,\mathrm{cm}$, $r = 1{,}0\,\mathrm{cm}$ og $m = 0{,}20\,\mathrm{g}$.

\begin{oppgaveliste}
\oppgave \fase{gjett}{Gjett}Kula er nøytral før den berører en plate. Hvorfor begynner den å sprette? Hvorfor spretter den tilbake igjen fra den negative plata?
\svar
\oppgave \fase{test}{Regn}Bruk ladningen kula får (se tallene under simuleringen). Regn ut feltstyrken, kraften og akselerasjonen. Kula må flytte seg $d - 2r = 6{,}0\,\mathrm{cm}$. Hvor lang tid bruker den på én tur? Sjekk med simuleringen.
\svar
\oppgave \fase{gjett}{Gjett}Ladningen kula får, er proporsjonal med feltstyrken. Hva skjer med kraften, tiden for én tur og middelstrømmen hvis du dobler spenningen? Test med 20~kV.
\svar
\oppgave \fase{forklar}{Forklar}Følger «ping-pong-kretsen» Ohms lov? Begrunn svaret ut fra hvordan strømmen avhenger av spenningen.
\svar
\oppgave \fase{test}{Undersøk}Gjør kula fire ganger så tung. Hva skjer med tiden for én tur og med strømmen? Hva skjer hvis du i stedet dobler radien til kula?
\svar
\oppgave \fase{test}{Regn}Kjør simuleringen i 10~s. Hvor mye ladning er ført over? Sammenlign med middelstrømmen ganget med tiden.
\svar
\oppgave \fase{utfordring}{Utfordring}Regn ut arbeidet feltet gjør på kula på én tur, og den kinetiske energien kula har når den treffer den andre plata. Hvor blir det av energien i et ekte forsøk?
\svar
\end{oppgaveliste}

\end{document}
